\section{Broader Impacts}
\label{sec:bi}

\textbf{MCU CTF Platform and Challenges.}
Existing CTF platforms, including pwn.college~\cite{PwnCollege}, OverTheWire~\cite{OverTheWire},
lack the architectural support and challenges for MCU systems. 
They are designed for the Intel x86/64 architecture only.
The PIs are developing and will host \ctfplatform,
a next generation and open CTF platform that supports challenges for MCU architectures. 
In particular, the PIs will develop control-flow hijacking and attack mitigation challenges for MCU systems, e.g., buffer overflow, heap exploitation.
\ul{Evaluation Plan:} The effectiveness of the platform will be measured based on the number of registered users and their level of engagement and activity. %Additionally, we will employ ad-hoc quantitative surveys and conduct qualitative interviews to gather valuable insights and feedback from users.

\textbf{Positively Impacting Curricula and Beyond.}
The PIs have created new courses, e.g., \emph{Software Vulnerabilities and Security}, \emph{Advanced Software Security}, \emph{Trusted Computing and Execution}, \emph{IoT Security},
using materials from both the literature and the PIs' research results.
The research results from this project and \ctfplatform will be naturally integrated into classroom teaching as demos, homework, projects, and exams.
Because systems and software security is difficult to learn, 
PI Zhao publicizes his course recordings, talks, and training recordings at the YouTube channel~\cite{youtubelist}, 
which has 400+ subscribers, 30,000+ views, and 3,000+ hours of watch time at the time of submitting this proposal.
To reach a wider audience, the PIs will continue to create research, education, and scientific outreach videos on this channel. 
\ul{Evaluation Plan:} The courses will be formally assessed by the success of classroom teaching with student evaluations. The number of subscribers and viewership statistics measure the impacts of the YouTube channel. 

\textbf{Economic/Societal Impacts and Technology Transfer}. 
MCU-powered embedded and IoT devices are omnipresent.
%For example, the ARM embedded and IoT ecosystem shipped 4.4 billion Cortex-M MCUs in a single quarter in 2020~\cite{Armnews2020}.
Meanwhile, %the worldwide spending on IoT reaches more than \$1 trillion by 2022~\cite{IDC}, and 
the IoT security spending amounts to \$1.2 billion in 2021~\cite{Statista}.
The research activities will fundamentally influence the way MCU security solutions are built and used.
As a result, the research not only saves costs in cyberattacks but also has broader societal impacts on energy, transportation, healthcare, public safety, and many other sectors. 
%To influence real-world systems, the PI will partner with Dr. Yong Ho Hwang (Corporate Vice President, Samsung Electronics) to evaluate and integrate the proposed algorithms and systems into the products of Samsung (the second largest electronics company in the world~\cite{top10elec}). 
Our prior work on software-based mitigation for cache side-channel attacks~\cite{cho2020smokebomb} has been evaluated on and integrated into Samsung Tizen OS~\cite{Tizen} for wearable devices, tablets, smart TVs, and in-vehicle infotainment systems.

\begin{wrapfigure}{r}{0.45\textwidth}  
     \centering
     \vspace{-.3cm}       
     \begin{subfigure}[b]{0.245\textwidth}
         \centering       
         \includegraphics[width=\textwidth]{MITREectf2023.png}         
         \caption{4th in MITRE eCTF'23}
         \label{fig:rank}
     \end{subfigure}
     \begin{subfigure}[b]{0.185\textwidth}
         \centering
         \includegraphics[width=\textwidth]{puhabi.jpg} 
         \caption{Puhabi}
         \label{fig:puhabi}
     \end{subfigure}     
     %\hfill
     \begin{subfigure}[b]{0.44\textwidth}
         \centering
         \includegraphics[width=\textwidth]{ectf.png}
         \caption{Team Cacti in eCTF'23 presentation}
         % (Zheyuan Ma, Gaoxiang Liu, Xi Tan, Md Armanuzzaman, Trevor Schupbach, Sagar Mohan, and Hiu Laam Chau)
         \label{fig:presentation}
     \end{subfigure}
     \vspace{-.35cm}      
        \caption{CTF Team and High School Advisee}
        \label{f:ectf}
        \vspace{-.35cm} 
\end{wrapfigure}
\textbf{Cybersecurity Training and CTF Teams with Underrepresented, Undergraduate, and High School Students.}
%The PI will continue to provide students with a safe infrastructure to perform cybersecurity research and develop MCU exploitation and defense skills.
Lead PI Zhao has founded the \emph{Cacti} hacking team with 300+ registered undergraduate and graduate students %(University at Buffalo and Rochester Institute of Technology) 
and over five high school students (Williamsville East high school, Buffalo; Brighton high school, Rochester; and Agartala high school, India).
%Most of registered students are undergraduates from Computer Science, Computer Engineering and Electrical Engineering departments.
%Team \emph{Cacti} competes in embedded and software CTFs. 
Team \emph{Cacti} placed 3rd in the MITRE embedded CTF (eCTF) 2019, 6th in 2020, 9th in 2021, 5th in 2022, 4th in 2023 (as shown in Figure~\ref{fig:rank} and Figure~\ref{fig:presentation}), and 4th in 2024.
%I will also recruit high school students from the Buffalo and Rochester metropolitan areas through the NSF-funded GenCyber Camps.
As shown in Figure~\ref{fig:puhabi}, PI Zhao's high school advisee Puhabi Chakraborti~(female) was conferred the Pradhan award by the Prime Minister Modi of India in 2022~\cite{puhabi2022} and received full scholarship to attend University at Buffalo in 2024.
\ul{Evaluation Plan:} 
The cybersecurity training impacts will be evaluated by student engagement, the number of CTFs team \emph{Cacti} participates, and the team's standings in the CTFs. 
The research experience with underrepresented, undergraduate, and high school students will be evaluated by publications.
